\section{High availability disabled}
\label{sec:ha:disabled}

The venue has a single switch that turns high availability off. With it off no
arbiter and no witness is started, and each component that would otherwise run
as a primary and secondary pair runs as one instance.

This configuration exists for development, and for no other purpose. A
development machine is one machine, and a second instance of every component
consumes processors and memory that the work being done has better uses for.
Many tests do not depend on high availability at all, and run faster and more
predictably without the components that provide it.

\textbf{It is not a supported production configuration and must never be
deployed as one.} With one instance of each singular component there is nothing
to fail over to. The only recovery available is to restart the process that died
and rebuild what it held.

\textbf{Most testing runs with high availability enabled.} The scenarios that
exercise the behaviour specified in this chapter run in that configuration,
because it is the configuration a production venue runs. This configuration is
tested for what is particular to it: that the components which exist only to
arbitrate are absent, and that the components which remain report their absence.

\textbf{Development is not restricted to this configuration.} High availability
can be switched on with every component running on one machine, and doing so
starts the arbiters, the witness and a second instance of each singular
component exactly as a production deployment would. A developer can therefore exercise the behaviour specified in
\cref{sec:ha:enabled} on one machine. Such a deployment cannot reproduce the loss
of a machine, because it has only one, so it exercises everything in that section
except \cref{sec:ha:enabled:machine_loss}.

\subsection{What the switch does and does not remove}
\label{sec:ha:disabled:scope}

Two kinds of multiplicity must not be confused.

\emph{Replicated peers} are components of which several instances run, all of
them live, none of them leading. The authentication services and the gateways
are replicated peers: a caller reaching any instance gets the same answer, and
losing one costs capacity rather than correctness.

\emph{Singular components} are those of which exactly one instance may act at a
time. The sequencer, the matching engine and the publisher are singular, and
that constraint is what the arbiter, the leader and follower roles, and the
leadership generation exist to enforce.

Turning high availability off removes the second instance of each singular
component, and the components whose only purpose is to arbitrate between them.
It says nothing about the replicated peers, which run in every configuration.

\subsection{What this configuration may and may not do}
\label{sec:ha:disabled:obligation}

This configuration is used for development and for running tests. Two
requirements follow from that, and they constrain the configuration in opposite
directions.

\textbf{This configuration may lose what a production venue must not.} It runs
one instance of each singular component, so the death of that instance can lose
an order the venue has taken, and can lose a session's numbering with it. That
follows from running one instance of everything, and it is accepted here because
nothing of value is being traded.

\textbf{It must not behave differently from production in ways a test cannot
see.} A test run against this configuration is evidence about the venue only if
the venue does the same things here that it does in production, and fails for the
same reasons. Where it does less, it must fail visibly rather than appear to
work.

The failure this rule prevents is a venue that accepts a member's orders,
acknowledges every one of them, forwards none of them anywhere, and logs no
error. A test that sent orders and received replies would pass against such a
venue, and the result would mean nothing.

\textbf{Only the primary instance of each singular component runs.} The gateways
send orders to the primary sequencer and to no other. No secondary instance has
any part in this configuration: there is no peer for it to follow and no arbiter to
make it lead, and if the primary's machine is lost the venue stops, because this
configuration does not cope with the loss of a machine.

\begin{requirement}{R-0163}{With high availability off, a secondary instance refuses to start}
  With high availability off, every secondary instance shall refuse to start:
  the secondary sequencer, the secondary matching engine and the secondary
  publisher. Each shall say that it refused because high availability is off. The
  replicated peers are not secondaries and are not affected: several
  authentication services and several gateways may run in this configuration as
  in any other.
  \rationale{A secondary started in this configuration would be a second instance that
  nothing sends orders to. The venue would accept members' logons and orders and
  answer none of them, which is the failure this section exists to prevent. The
  arbiter and the witness refuse to start in this configuration for the same
  reason.}
  \covers{ha_test:46}
  \covers{ha_test:47}
\end{requirement}

\subsection{Recovery in general}
\label{sec:ha:disabled:recovery}

Two requirements hold for every component in this section, and this subsection
states them once rather than repeating them for each component. Both concern the
period between a component starting and that component being ready. That period
is short, every restart passes through it, and a second failure during it does
the most damage.

Every requirement here assumes that a component can read what survived the
failure. A log truncated part-way through a record, a checkpoint that fails its
checksum, or a file that is absent leaves a component with nothing to recover
from. This configuration provides no second instance holding an undamaged copy.

\textbf{This configuration does not protect against damaged or missing durable
state.} Where high availability is enabled, the instance whose copy is
undamaged takes over, and that is the venue's answer to damaged state. The
disabled configuration has no equivalent, so it should be run only where losing
what it holds is acceptable.

\begin{gap}{}
  Damaged durable state is not detected as such. A component reads what it can
  and proceeds, rather than establishing that what it read was whole, so a
  truncated log or a partial checkpoint produces a component that holds the wrong
  orders and reports itself as current. A member is then answered from that
  state: an order it placed may be unknown, and one it cancelled may be live.
\end{gap}

\begin{requirement}{R-0121}{A component takes the cost of first touching its memory before it reports itself ready}
  Where a component maps memory it will use while trading, it shall touch all of
  that memory before reporting itself ready, so that the delay of first use falls
  before it is trading rather than during.
  \rationale{Mapping a file reserves addresses and reads nothing. The first time a component
  uses each page, the processor interrupts and the kernel fetches that page. The
  delay is small and happens once per page, and a region of a few hundred megabytes
  holds tens of thousands of pages. Where the component does not touch them first,
  those delays fall on whatever code reaches each page first, which is the order
  path.

  Touching the memory before reporting ready means that a component which has said
  it can trade has already taken those delays. Touching it afterwards moves the
  cost onto the first members to send anything, which is the worst moment to take
  it and the hardest to attribute afterwards.}
  \uncovered{no scenario measures a component's first orders against its later ones}
\end{requirement}

\begin{requirement}{R-0122}{A report a member receives twice is marked as a repeat}
  Where a recovery repeats work whose execution reports had already been sent,
  the reports sent again shall be marked so that the member can tell a repeat
  from a second event.
  \rationale{A component that dies part-way through handling an order holds no record of
  whether it sent the report for that order. Sending the report again is the right
  choice. A member that never receives a report cannot recover from that on its
  own. A member that receives one twice can recover, provided it can distinguish
  the two.

  Where the venue does not mark the repeat, a member cannot distinguish a repeated
  acceptance from a second order accepted. The venue would then report an
  acceptance that did not happen, while correcting one that did.}
  \covers{ha_test:53, ha_test:99}
\end{requirement}

\begin{gap}{}
  Scenario 53 covers the repeat that a catch-up produces, which is the common case
  and the one a member meets. An engine applies the venue's record and marks every
  report it sends, because it cannot distinguish a record an earlier engine
  reported from one no engine ever saw.

  It does not cover the narrower case that the requirement also states: a component
  that dies between sending a report and recording that it sent it, so that the
  same work is repeated from a record that is one step behind. No scenario kills a
  component in that window, and producing the window deliberately is harder than
  the other scenarios, because it is a few instructions wide.
\end{gap}

\begin{requirement}{R-0031}{A component does not act on partially recovered state}
  A component shall not begin work until its recovery is complete. No sequence
  number shall be assigned before the last one issued is known, and no order
  shall be accepted before the set of open orders is current.
  \rationale{Acting on partial state produces results that look ordinary and are wrong. Such
  results are the hardest of the venue's failures to detect, and the hardest to
  reverse afterwards.}
  \uncovered{no scenario interrupts a recovery}
\end{requirement}

\begin{requirement}{R-0032}{An interrupted recovery can be repeated without applying anything twice}
  Where a component fails while recovering, the recovery shall be able to run
  again from the same starting point and reach the same result.
  \rationale{Without this requirement, a failure during recovery corrupts state rather than
  delaying the component, and a component that fails repeatedly leaves worse state
  after each attempt. This requirement is also what makes a bound on the number of
  attempts safe, because the bound protects the venue only where the attempts
  beneath it did no harm.}
  \uncovered{no scenario interrupts a recovery}
\end{requirement}

\subsection{The authentication service}
\label{sec:ha:disabled:auth}

Several authentication services run at once, and none of them leads. A member
never contacts one directly. A gateway contacts an authentication service on the
member's behalf at logon, and reaches whichever instance answers.

Losing one instance reduces the capacity available and does not affect
correctness. Two instances holding different answers to the same question does
affect correctness, and most of this subsection concerns that condition.

\subsubsection{What state it holds}
\label{sec:ha:disabled:auth:state}

An instance holds the credential set: which comp ids may trade, the material each
is authenticated against, and the instructions each gave when it was admitted.
This set is a
copy rather than the original. The database holds the original, and the
administration service is the only component that writes to it. That service
changes credentials and distributes each change to every authentication service,
and \cref{sec:ha:disabled:admin} describes it.

\subsubsection{Where the copy that outlives it lives}
\label{sec:ha:disabled:auth:durable}

The database holds the copy, and an instance reaches it in two steps rather than
one. An instance loads an export of the database when it starts, and the
administration service afterwards sends it each change as it makes it. The
instance holds no connection to the database, so it can neither check an answer
against the original nor establish how far behind it has fallen.

\textbf{Only the administration service can establish whether an instance is
current}, because only that service holds a record of what it sent and to which
instance.

\begin{gap}{}
  A change is reported as having succeeded when at least one instance applied
  it. A partial fan-out is therefore indistinguishable from a complete one, and
  nobody learns that an instance was missed.
\end{gap}

\begin{gap}{}
  Nothing brings a running instance up to date. Having missed a change, it stays
  wrong until it is restarted --- and restarting reloads the export, which may
  itself predate the change, so restarting is not a reliable remedy either.
\end{gap}

\subsubsection{How a restarted instance becomes current}
\label{sec:ha:disabled:auth:recovery}

The instance reloads the export and resumes answering. It has one such operation
rather than two, as the gateway in \cref{sec:ha:disabled:gateway} does, because
an authentication service holds no role and is never promoted. A restarted
instance and a newly started one therefore do the same thing.

The instance does not ask for the changes it missed, and no component offers them
to it.

\subsubsection{What the member is told}
\label{sec:ha:disabled:auth:member}

An instance can be wrong in two ways, and the two are not equally serious, which
is why the requirements below treat them differently. An instance that has missed
an addition refuses a member that should be admitted. The member sees that
refusal at once, and the venue can recover from it. An instance that has missed a
withdrawal admits a member that should be refused. Nothing reports that
admission, and it persists until somebody corrects the instance.

\begin{requirement}{R-0042}{A credential change is applied by every instance, or is reported as incomplete}
  A change shall not be reported as having succeeded unless every instance
  applied it. Where an instance did not, the change shall be reported as
  incomplete and shall name the instance.
  \rationale{Reporting success when one instance of several applied the change leaves a
  partial distribution looking the same as a complete one, so nobody can act on it.
  Every remedy below depends on the venue knowing that the condition exists.}
  \uncovered{no scenario makes a credential change while an instance is down}
\end{requirement}

\begin{requirement}{R-0043}{An instance that has missed a change is brought current without being restarted}
  Where an instance has missed a change, it shall be brought up to date while
  running.
  \rationale{Restarting an instance is not a remedy, because the instance reloads an export of
  the database that may itself be older than the change it missed. Nothing else
  brings an instance up to date, so without this requirement an instance stays
  stale indefinitely.}
  \uncovered{no scenario makes an instance miss a change, and nothing would bring it current}
\end{requirement}

\begin{requirement}{R-0044}{A member is admitted by every instance or by none}
  A member's logon shall succeed or be refused alike whichever instance answers
  it.
  \rationale{Intermittent admission is worse for a member than consistent refusal. The logon
  succeeds, then it fails, and nothing distinguishes the two attempts, so the
  member can neither report the fault clearly nor reproduce it.}
  \uncovered{no scenario authenticates the same member against both instances}
\end{requirement}

\begin{requirement}{R-0045}{An instance can say how current its credential set is}
  An instance shall be able to report the point up to which it has applied
  changes.
  \rationale{Where an instance cannot report this, no component can establish it. The instance
  has no view of the database, and the administration service cannot distinguish an
  instance that missed a change from one that never received it.}
  \uncovered{no scenario asks an instance how current its credentials are, and nothing answers}
\end{requirement}

\begin{requirement}{R-0046}{A change made while an instance is unavailable is applied when it returns}
  A credential change shall be accepted while an instance is unavailable, that
  instance shall be recorded as not current, and the change shall be applied to
  it when it returns.
  \rationale{Refusing the change instead would stop the venue withdrawing access during an
  outage, which is when withdrawing access may matter most. The cost is an instance
  that is behind for a bounded period. That cost is acceptable because
  \reqref{R-0042} makes the venue aware of the condition.}
  \uncovered{no scenario makes a credential change while an instance is down}
\end{requirement}

\subsection{The administration service}
\label{sec:ha:disabled:admin}

The administration service is the only component permitted to change credential
state. Members never contact it. Authentication services receive what it sends
them.

\textbf{One instance runs in both configurations, by decision.} The service is
not on the path an order takes, so its absence stops no member trading or logging
on. Credentials change rarely, so the venue tolerates its absence for as long as
a restart takes. A second instance would need arbitration of its own to remain
the single writer, which is a large mechanism for a component that is idle for
most of a trading day.

Because the administration service has only one instance, any state it must
retain after a restart must be stored durably outside its running process.

\subsubsection{What state it holds}
\label{sec:ha:disabled:admin:state}

\textbf{The credential state.} The service writes this state, and the database
holds the original.

\textbf{Which authentication services have applied which change.} No other
component holds this. An authentication service holds no record of a change it
never received, and the database records what the credentials are rather than
which instance has been told.

\subsubsection{Where the copy that outlives it lives}
\label{sec:ha:disabled:admin:durable}

The database holds the credential state. Nothing holds the record of which
instances have applied which change.

\begin{gap}{}
  Which instances have applied a change is held only in the running process. An
  administration service that restarts between making a change and distributing
  it to every instance cannot afterwards establish which instances were missed.
\end{gap}

\subsubsection{How a restarted instance becomes current}
\label{sec:ha:disabled:admin:recovery}

The service reads the credential state from the database, and that state is
complete. It cannot recover the record of what it had distributed, because
nothing wrote that record down.

The service has one such operation rather than two, as the gateway in
\cref{sec:ha:disabled:gateway} and the authentication service in
\cref{sec:ha:disabled:auth} do. An administration service is never promoted,
because the venue never runs a second one to promote.

\subsubsection{What the member is told}
\label{sec:ha:disabled:admin:member}

Nothing directly. An authentication service answers a member's logon, and the
requirements in this subsection exist so that the answer does not depend on an
administration service running at the time.

\begin{requirement}{R-0096}{A credential change is recorded before it is distributed}
  A change shall be written to the source of truth before it is sent to any
  instance, so that a failure between the two loses the distribution rather than
  the change.
  \rationale{The next instance to reload from the database undoes a change that was
  distributed but not recorded, and nothing reports that. A change that was
  recorded but not distributed appears as a difference between instances, and the
  venue can correct it.}
  \uncovered{no scenario interrupts a credential change}
\end{requirement}

\begin{requirement}{R-0097}{The record of which instances are behind outlives the service that made it}
  Where an instance did not receive a change, that shall be recorded durably, and
  shall be acted on by whichever administration service is running when the
  instance next becomes reachable.
  \rationale{\reqref{R-0046} accepts a change while an instance is unavailable, on the basis
  that the venue applies it later. Where only a running process holds the
  obligation to apply that change, the obligation is lost whenever the
  administration service restarts before the instance returns. The credential state
  is then permanently inconsistent, and nothing records that it is.}
  \uncovered{no scenario restarts the administration service}
\end{requirement}

\begin{requirement}{R-0112}{The venue trades while the store it depends on is unreachable}
  No component shall require a connection to the durable store in order to admit
  a member, accept an order, or answer one. What it needs from the store shall be
  read when it starts and supplied to it thereafter.
  \rationale{Where the venue needs the store to be reachable in order to work, every outage of
  the store becomes a trading outage, and the venue takes on a failure whose
  likelihood has nothing to do with anything the venue does. Admission already
  works this way: an authentication service reads an export when it starts and is
  sent each change afterwards. Orders that outlive the day are to work the same
  way.}
  \uncovered{no scenario runs the venue with the store unreachable}
\end{requirement}

\begin{requirement}{R-0098}{Members trade and log on while no administration service is running}
  The absence of an administration service shall not prevent a member logging on,
  placing orders, or being answered.
  \rationale{The administration service writes credential state, and answers nothing when a
  member is admitted. Making a member's logon depend on it would place a component
  that is idle for most of a trading day on the path of every session that starts.}
  \uncovered{no scenario runs the venue without an administration service}
\end{requirement}

\subsection{The gateway holding the session}
\label{sec:ha:disabled:gateway}

The gateway carrying a member's session dies and is restarted. The member's
connection ends and it reconnects.

\subsubsection{What state it holds}
\label{sec:ha:disabled:gateway:state}

The gateway holds the state of the session: how far the session's numbering has
reached in each direction, and which of the outbound numbers carried execution
reports rather than session-level traffic. This is the only state that matters
when the gateway dies.

While it is running, the gateway also holds the member's identity and the
instructions the member gave when it was admitted --- whether its resting orders
should be cancelled if its connection ends, and how long the venue should wait
before doing so --- any authentication exchange part-completed, any
resend being served, and any reports held back while the numbering is being
established. None of these needs to survive the gateway. The member's identity is
established again at the next logon, and the venue abandons the rest, which the
member asks for again.

The gateway holds no order state. The matching engine holds which orders are
open, and the sequencer holds the order they were taken in. Where the gateway
must act on a member's behalf, such as cancelling what the member has resting
when its connection ends, it reports the disconnection and those instructions to
the component that owns the orders. It keeps no view of those orders to act on
itself. The instructions therefore belong to the session rather than to the
gateway.

Those instructions reach the gateway when the member is admitted, and a component
needs them later, possibly a different instance, and possibly when the member is
no longer connected to supply them again. The requirement \reqref{R-0103} states
what must be true wherever they are held, and this chapter does not say where
that is. \texttt{docs/availability/session\_binding.md} makes that choice and
gives the reason.

Storing the instructions where a gateway can retrieve them handles a restart
only after the member reconnects and binds to a gateway. It does not handle a
gateway that dies while holding live sessions: a member who never reconnects
will not trigger another bind. In that case, the component that owns the orders
must carry out the member's instruction when the gateway dies (\reqref{R-0113}).

\subsubsection{Where the copy that outlives it lives}
\label{sec:ha:disabled:gateway:durable}

The sequencer holds the numbering. A gateway reports a session's position as it
advances, and the sequencer records it, because the sequencer is the only
component that every gateway instance talks to. The sequencer's write-ahead log
holds the execution reports, each stamped with the session it belongs to.

An order that has not been committed has no copy anywhere in the venue. Between
the member sending it and the sequencer committing it the only record inside the
venue is the gateway's, and when the gateway is what died, nothing can recover
it.

This follows from a deliberate decision. A durable record kept by the gateway
would add a write to the hot path, in front of the write-ahead log that already
exists, and would double the cost of every order to protect a window of
microseconds. The member that sent the order also holds a copy of it outside the
venue. The venue therefore does not attempt to recover an uncommitted order. It
guarantees that such an order never takes effect afterwards, that the member can
establish this, and that the member can send the order again.

\begin{requirement}{R-0008}{The most recent messages to each member remain recoverable}
  For each comp id, the venue shall keep the most recent application messages it
  has sent to that comp id during the current trading day, up to 100,000 messages.
  The count shall be one figure that applies to every comp id. A resend request
  for a message older than those kept shall be answered with a FIX gap fill, a
  SequenceReset with GapFillFlag set, covering the messages that are no longer
  kept. No message from an earlier trading day can be asked for.
  \rationale{Recovery that depends on a record still being present is only as good as the
  retention behind that record. A stated count gives a member a figure to rely on,
  where retention that follows from how much storage happens to be free gives it
  nothing. A fixed count for every member, with a gap fill beyond it, is the
  established practice at the London Stock Exchange and at Turquoise, so members
  already know how to work with it. Because the count is in messages rather than
  time, a member that is sent many messages has a shorter window in time than a
  quiet one. At no more than two messages per order, 100,000 messages last about
  eight minutes for a member sending 100 orders a second, and about fifty seconds
  for one sending 1,000.}
  \uncovered{no scenario measures how far back messages remain available}
\end{requirement}

\begin{gap}{BUG-0048, BUG-0113}
  Nothing keeps messages for each comp id separately. A resend is served from the write-ahead log, which discards nothing,
  so a message stays available for as long as the venue runs and the disk holds
  out, and no gap fill is ever sent for a message that is too old.
\end{gap}

\subsubsection{How a restarted instance becomes current}
\label{sec:ha:disabled:gateway:recovery}

The member reconnects, the session binds to the restarted gateway, and the
sequencer supplies the numbering it holds. What the member receives next depends
on how far each of its orders had got before the gateway died.

The gateway has one such operation where other components have two, because a
gateway is never promoted and holds no role. A member returning to the instance
that restarted, and a member arriving at a different instance, are the same event
from the venue's side: a session binds to an instance that does not hold its
state, and that instance asks the sequencer for it. Nothing in that sequence
depends on how many gateways are running, so gateways need no arbiter and running
several of them costs little.

\textbf{An order the gateway had received but not forwarded} was never
committed.

\begin{requirement}{R-0001}{An order that was not forwarded never takes effect}
  The order shall have no effect on the venue. No execution report shall be
  produced for it, and the venue shall hold no order open on its behalf.
  \rationale{The venue had not taken the order. Acting on it later would execute an
  instruction that the member has been given no reason to believe is live.}
  \uncovered{no scenario restarts a gateway}
\end{requirement}

\textbf{An order the sequencer had committed} is the venue's responsibility from
that moment.

\begin{requirement}{R-0004}{A committed order is processed}
  The order shall be processed by the matching engine, whether or not the
  gateway that received it is running when the engine reaches it.
  \rationale{Committing the order is the venue taking it. What happens after that point must
  not depend on which component carried the order.}
  \uncovered{scenario 23 covers a gateway that is not restarted}
\end{requirement}

\textbf{A report that had been given an outbound sequence number} but did not
reach the member leaves a gap in the member's numbering.

\begin{requirement}{R-0007}{A numbered report that was not received is recoverable}
  The member shall detect the gap in its inbound numbering and ask for the
  missing range, and the venue shall send again the execution reports that
  occupied those numbers, filling the numbers that carried session traffic
  rather than reports.
  \rationale{This is the FIX resend, used as the protocol intends. The venue supplies the
  reports it can reproduce, and accounts for the numbers it cannot.}
  \covers{ha_test:22}
\end{requirement}

\subsubsection{What the member is told}
\label{sec:ha:disabled:gateway:member}

\begin{requirement}{R-0103}{The instructions a member gave at logon outlive the component that received them}
  The instructions a member gave about what should happen to its orders when its
  connection ends shall be available to whichever component must act on them,
  whether or not the component that first received them is still running.
  \rationale{Those instructions arrive once, when the session is admitted, and a component
  acts on them later, possibly a different instance, and possibly when the member
  is no longer present to supply them again. Instructions held only by the
  component that received them are lost in exactly the failure they guard against.}
  \covers{ha_test:107}
\end{requirement}

\begin{requirement}{R-0113}{A member's instruction is carried out when the component holding it is the one that failed}
  Where a member has asked for its orders to be cancelled if its connection
  ends, and the connection ends because the component holding that session
  failed, the instruction shall still be carried out.
  \rationale{A component acts on the instruction by reporting the disconnection to whatever
  holds the orders, and a component that has died reports nothing. The death of a
  gateway ends every session it holds, and is therefore the one failure after which
  nothing acts on what those sessions asked for. The member cannot detect this,
  because from the member's side the connection ended in the ordinary way.}
  \covers{ha_test:107}
\end{requirement}

\begin{requirement}{R-0002}{A member can establish that an order was not taken}
  On reconnecting, a member shall be able to determine that the venue holds no
  order for a \texttt{ClOrdID} it had sent.
  \rationale{The member cannot distinguish this case from a lost reply. Every
  venue provides an order status enquiry for exactly this reason.}
  \uncovered{no scenario asks the venue what it holds for a member, because nothing answers}
\end{requirement}

\begin{gap}{BUG-0089}
  The venue answers no order status enquiry of any kind, so a member has no way
  to ask what is held for it. Recovery rests entirely on what the venue chooses
  to send.
\end{gap}

\begin{requirement}{R-0003}{An order that was not taken may be resubmitted unchanged}
  A member may send the order again under the same \texttt{ClOrdID}, and the
  venue shall process it as a new order.
  \rationale{Without this requirement, the member must choose between abandoning an order and
  risking a duplicate. Resubmitting under the original identifier removes that
  choice.}
  \uncovered{no scenario restarts a gateway}
\end{requirement}

\begin{requirement}{R-0006}{A resubmitted committed order is refused as a duplicate}
  Where a member sends again an order the venue had already committed, under the
  same \texttt{ClOrdID}, the venue shall refuse it as a duplicate and shall not
  create a second order.
  \rationale{This requirement is what makes \reqref{R-0003} safe. A member cannot tell the two
  cases apart, and relies on the venue to distinguish them.}
  \covers{ha_test:98}
\end{requirement}

\begin{requirement}{R-0010}{A duplicate is refused whatever has become of the first order}
  The refusal required by \reqref{R-0006} shall not depend on the order having
  reached the matching engine, nor on the engine that received it still running.
  \rationale{The sequencer is the component that commits an order, so only a record of what
  has been committed can answer this question at the moment it is asked. The
  matching engine does not hold an order that has been committed and not yet
  applied, and an engine that has restarted holds no record of an earlier
  submission. In either case a member following \reqref{R-0003} would be given a
  second live order.}
  \covers{ha_test:96, ha_test:98}
\end{requirement}

\begin{requirement}{R-0119}{An order identifier used today is refused for the rest of the day}
  Where a member sends an order under an identifier it has already used on that
  session that trading day, the venue shall refuse it, whatever became of the
  first order and whatever has restarted since.
  \rationale{Two live orders under one identifier make a cancel request ambiguous, and a check
  against the orders currently open would catch that much. The reason to go further,
  and refuse an identifier whose first order has finished, concerns the record the
  venue publishes. Consumers of that record identify an order by this value, and a
  consumer such as a member's back office, a reconciliation process or an auditor
  cannot ask the venue anything. Two different orders under one identifier in that
  record cannot be read correctly.

  This requirement also makes \reqref{R-0003} safe. This document requires
  resubmission under the original identifier as a recovery path, for a member that
  does not know whether the venue took the order. The check that distinguishes a
  resubmission from a new order is therefore something a member following this
  specification depends on. If that check does not survive a restart, a member
  following the specification is given a second live order at the moment it is
  least able to notice.

  The trading day bounds the set of identifiers. The venue adds identifiers through
  a trading day and discards the set at the end of it, which keeps the set finite.}
  \covers{ha_test:96, ha_test:98}
\end{requirement}

\begin{requirement}{R-0005}{A report produced for an unbound session is delivered on reconnection}
  An execution report produced while the session had no connection shall be
  delivered when that session next binds to a gateway, without the member having
  to ask for it.
  \rationale{The venue allocated no outbound sequence number for such a report, so it leaves
  no gap in the member's numbering. The member has nothing to notice and nothing to
  request, and a resend cannot reach it.}
  \covers{ha_test:70}
\end{requirement}

\begin{gap}{BUG-0046}
  The FIX gateway asks for these reports once a session is established, and the
  leading sequencer sends them. The binary order gateway does not ask, so a
  member using the binary protocol is not sent the reports produced while it had
  no connection.
\end{gap}

\begin{requirement}{R-0009}{A gap that cannot be filled is reported as such}
  Where the venue cannot supply a report the member has asked for, it shall say
  so rather than fill the number silently.
  \rationale{A gap filled without explanation leaves the member believing its numbering is
  complete when it has lost a report it will never see.}
  \uncovered{no scenario restarts a gateway}
\end{requirement}

\subsection{The sequencer}
\label{sec:ha:disabled:sequencer}

The sequencer dies and is restarted. This breaks no member's connection, because
gateways hold the sessions and the gateways stay up. What stops is the venue's
ability to take an order.

\subsubsection{What state it holds}
\label{sec:ha:disabled:sequencer:state}

\textbf{The ordering.} The sequencer holds the number it will assign next, and
the write-ahead log that records what it has assigned. This log is the venue's
record of what it has taken. The rest of the state in this subsection depends on
it.

\textbf{The leadership generation.} The sequencer holds the generation its work
is stamped with, so that a receiver can distinguish a message from a superseded
instance from a message from the current one.

\textbf{Each session's numbering.} The sequencer holds how far every session has
reached in each direction, and which of its outbound numbers carried execution
reports. It holds this on the gateways' behalf rather than for its own use, as
\cref{sec:ha:disabled:gateway:durable} describes, because it is the only
component that every gateway instance talks to.

\textbf{Where an execution report should go.} The sequencer holds the route from
a committed order back to the session that placed it.

\textbf{What it holds while no matching engine leads.} The sequencer holds the
orders it has accepted but not yet sent to an engine while no matching engine
is the leader. It tracks how many orders are waiting and how long no engine has
been leader. It sends the waiting orders when an engine becomes leader.

\subsubsection{Where the copy that outlives it lives}
\label{sec:ha:disabled:sequencer:durable}

\textbf{The ordering and the generation survive the process.} The log is on
disk, and the sequencer takes the next number from the log's last record. The
generation is in a small file beside the log. The sequencer writes that file by
writing a temporary file and renaming it over the real one, so a reader sees
either the old value or the new one, whatever moment the process died at.

\textbf{The other three have no durable copy.} The running process holds them
and nothing else does. A restart treats each of the three differently.

The sequencer \emph{relearns} session numbering rather than recovering it.
Gateways report every session's position on a timer, so a restarted sequencer is
told again, within one reporting interval, about every session bound to a gateway
at that moment. No gateway reports a session that has no live connection, and
that is the session about to reconnect and ask for its numbering.

\begin{gap}{}
  A session that was not connected when the sequencer restarted is forgotten
  entirely. Nothing on disk holds its numbering and no gateway is reporting it.
\end{gap}

The sequencer does not rebuild the route from an order to a session, and that is
deliberate. Rebuilding it by replaying the log would fill it with entries that
nothing removes, and the log holds a trading day of them.

\begin{gap}{}
  An execution report arriving for an order committed before the restart cannot
  be placed, because the route it needed was not kept and is not rebuilt.
\end{gap}

What the sequencer holds while no engine leads is kept in memory. The orders
themselves are in the log, and a restarted sequencer holds no record that any of
them is still to be answered. It does not need one. A matching engine presents the position it
has reached, and the sequencer serves everything after that position from the
log, so which records an engine receives is determined by the position the engine
presents rather than by anything the sequencer remembers.

\subsubsection{How a restarted instance becomes current}
\label{sec:ha:disabled:sequencer:recovery}

The sequencer reopens the log, takes the next number from the last record it
finds, and reads the generation from the file beside the log. Other components
supply what it cannot recover: gateways report the sessions they hold, and a
matching engine connects.

What is outstanding at the moment the sequencer dies falls into two cases, and
the venue treats them differently.

\textbf{An order committed and not yet forwarded to the matching engine} belongs
to the venue, and \reqref{R-0013} covers it. The engine reports how far it has
reached, and the sequencer sends the records that follow.

\textbf{An execution report produced but not yet sent onward} is the second
case. The report survives, because the sequencer wrote it to the log stamped with
the session it belongs to. The route the sequencer was going to use does not
survive.

\begin{requirement}{R-0025}{A report can be addressed to its session after a restart}
  An execution report the venue has produced shall be deliverable to the session
  that placed the order, after a restart of any component between the two, and
  the session it belongs to shall be determinable from what survived that
  restart.
  \rationale{The record of the report already names its session. A route held only in a
  running process is lost at the moment it is needed, and no other component can
  supply it.}
  \uncovered{no scenario restarts a sequencer with reports undelivered}
\end{requirement}

\textbf{A promotion is a different operation, and it recovers more.} Gateways
report to both instances, so a follower has been receiving each session report
all along, and has been receiving every log record as it was written. A follower
therefore already holds what a restarted instance must be told, including the
numbering of sessions that are bound to no gateway, which are the sessions no
gateway will report.

The gateway in \cref{sec:ha:disabled:gateway} has one such operation, because a
gateway is never promoted. The sequencer has two. They recover different amounts,
and the one that recovers less is the one the venue uses when it has no second
instance to fail over to.

\subsubsection{What the member is told}
\label{sec:ha:disabled:sequencer:member}

\begin{requirement}{R-0026}{A logon completes or is refused within a bounded time}
  A member logging on shall be given either an established session or a refusal,
  within a stated time, whatever the venue has restarted meanwhile. It shall not
  be left authenticated and unestablished.
  \rationale{A session that is neither established nor refused leaves the member unable to
  trade and with nothing to act on. The member cannot proceed, and it has been
  given no reason to reconnect.}
  \uncovered{no scenario restarts a sequencer during a logon}
\end{requirement}

\begin{requirement}{R-0027}{An interrupted resend leaves the member able to act}
  Where the venue cannot finish supplying a range of messages it has begun, it
  shall say so. A member shall not be left with an unfilled gap and no
  indication that nothing further is coming.
  \rationale{From the member's side, waiting for records that will never arrive looks the same
  as waiting for records that are slow. Only the venue can distinguish the two
  cases, so only the venue can end the wait.}
  \uncovered{no scenario restarts a sequencer during a resend}
\end{requirement}

\begin{requirement}{R-0011}{A session's numbering survives a restart of whatever remembers it}
  A member reconnecting after the venue has restarted the component that holds
  its numbering shall be resumed where it left off, and not at the beginning.
  \rationale{A member whose numbering restarts sees a break it cannot reconcile, and the venue
  cannot tell that member which of its orders the venue still holds.}
  \uncovered{no scenario restarts a sequencer with a session unbound}
\end{requirement}

\begin{requirement}{R-0012}{A member is never resumed at a numbering the venue cannot vouch for}
  Where the venue has no record of a session's numbering, it shall say so rather
  than resume that session at the beginning.
  \rationale{A member whose own store has also restarted cannot distinguish a session resumed
  at the beginning from a session that has never traded. Both sides then agree on a
  numbering that neither has any basis for, neither side sees a gap, and the
  member's numbering is reset while its orders are still open. A member can recover
  from being refused. A member cannot recover from being agreed with wrongly.}
  \uncovered{no scenario restarts a sequencer with a session unbound}
\end{requirement}

\begin{gap}{}
  A member whose session the venue has forgotten is resumed at the beginning and
  is not told.
\end{gap}

\begin{requirement}{R-0013}{An order committed before a sequencer restart still receives its execution report}
  An order that the sequencer had committed before it died shall be processed,
  and an execution report shall be produced for it and be deliverable to the
  session that placed the order.
  \rationale{Committing the order is the venue taking it. Restarting the component that took
  the order does not return the order to the member.}
  \uncovered{no scenario restarts a sequencer with orders committed and unanswered}
\end{requirement}

\begin{requirement}{R-0014}{A sequence number is never issued twice}
  The venue shall not assign a sequence number it has assigned before, whatever
  has restarted since.
  \rationale{The sequence number is the venue's name for an order. The venue cannot answer two
  orders that share one number, and cannot distinguish the second from a
  retransmission of the first.}
  \uncovered{scenario 14 covers a restart with a peer to synchronise from}
\end{requirement}

\subsection{The matching engine}
\label{sec:ha:disabled:me}

The matching engine dies and is restarted. No member's connection is broken, and
the venue can still take orders. What stops is the processing of those orders.

\subsubsection{What state it holds}
\label{sec:ha:disabled:me:state}

\textbf{The set of open orders.} The engine holds which orders the venue has
accepted and has not yet cancelled or rejected, and the session each order
belongs to. The engine uses this set to recognise an order identifier that is
already in use, and to answer a cancel request. This set is the only copy of it.
The sequencer holds a record of which orders were taken and in what order, which
is a different thing, and converting that record into this set is expensive.

\textbf{The point the set has reached.} The engine holds the sequence number up
to which the set reflects what the sequencer has committed. Without that number
the engine cannot establish which committed orders are missing from the set.

\subsubsection{Where the copy that outlives it lives}
\label{sec:ha:disabled:me:durable}

In a memory-mapped region of fixed-size records, one for each open order, held
beside the engine and separate from the structures it works with. An order is
written to it as it is accepted and taken out of it as it is cancelled, and it is
read back when the engine starts. See
\texttt{docs/durability/open\_order\_checkpoint.md}.

The region survives the death of the process. It does not survive the death of
the machine, and this configuration is specified to give no more than that.
Where high availability is enabled, the answer to a lost machine is the other
machine.

\textbf{Reading the region is what a restart costs.} Scenario 50 recovers a
thousand open orders from a region of 248 MB in 56 milliseconds, where the pages
the dead process wrote are still held in memory. The same walk takes 457
milliseconds where those pages have to be read from the disk. Both figures follow
from the size of the region rather than from the number of orders open in it.

\textbf{The region holds state, and the sequencer's log holds changes.}
Rebuilding the set of open orders from the log instead means replaying the log
from the start of the trading day. The size of the region follows from the number
of orders open at one moment. The cost of a replay follows from the number of
orders the day has seen, which is far larger and grows all day, so a replay at a
full day's volume takes tens of seconds where reading the region takes
milliseconds.

The venue does replay the log in one case, which \reqref{R-0123} states: where
the region cannot be used at all. It replays the log from the first record, and
where the log reaches back to the first command of the trading day, the orders it
rebuilds are exactly the ones the venue held, so it keeps them and carries on.
Where the log does not reach back that far, it halts trading.

\begin{gap}{BUG-0048}
  Nothing bounds how far back the log is kept. The venue discards no records at
  all, so the log grows for as long as the venue runs, and what it retains is
  anchored to no position any component may present.
\end{gap}

\textbf{A checkpoint records the set itself}: the orders the component holds
open, and the sequence number those orders are current to. The region is that
checkpoint. It is the copy that outlives the process, and the log supplies only
what has happened since.

\subsubsection{How a restarted instance becomes current}
\label{sec:ha:disabled:me:recovery}

The engine reads the checkpoint and rebuilds the set as it stood at the sequence
number the checkpoint records. It then asks the sequencer for everything
committed after that number. It applies what comes back without issuing execution
reports, because those reports were issued the first time. It begins accepting
orders once it has caught up.

\textbf{A promoted instance performs the same operation.} A follower that is
promoted already holds a replica of the set and the sequence number that replica
reached. It asks for everything after that number, applies the result silently,
and begins accepting orders. The two cases differ only in where the starting
number comes from: a file after a restart, and a replica after a promotion.

Two things follow. The configuration with high availability enabled needs no
separate promotion path, because promotion is this same recovery starting from a
number learned another way. The mechanism also runs on every restart, so it is
exercised continually rather than only when a test produces a failover.


\subsubsection{What the member is told}
\label{sec:ha:disabled:me:member}

\begin{requirement}{R-0028}{Every accepted order is reported to its member, once}
  An order the matching engine accepted shall have an execution report reach the
  member that placed it, whatever restarted between the acceptance and the
  delivery, and the member shall not be led to believe the acceptance happened
  twice.
  \rationale{Catching up after a restart applies orders without issuing reports,
  on the ground that the reports were issued the first time. An order accepted
  whose report never left the engine is exactly the case where that ground does
  not hold, so something must be able to tell the two apart.}
  \uncovered{no scenario restarts the engine between acceptance and delivery}
\end{requirement}

\begin{requirement}{R-0029}{A cancel request is always answered}
  A request to cancel shall receive either a report saying the order was
  cancelled or a rejection saying why, whatever restarted while the request was
  in flight.
  \rationale{A member that does not know whether its cancel took effect must assume the order
  is live, and it cannot hedge in either direction. A member can act on a
  rejection. It cannot act on silence.}
  \uncovered{no scenario restarts the engine with a cancel in flight}
\end{requirement}

\begin{requirement}{R-0030}{With no matching engine, the venue refuses rather than acknowledges}
  Where no matching engine exists to process orders, the venue shall refuse
  orders and tell members why, rather than accept orders it cannot process.
  \rationale{An order acknowledged and never processed is worse for a member
  than one refused: it believes the order live, may have hedged against it, and
  has nothing to cancel.}
  \uncovered{scenario 42 covers refusal, and not a restart}
\end{requirement}

\begin{requirement}{R-0018}{An open order survives a restart of the matching engine}
  An order the venue held open before the restart shall be held open afterwards,
  on the same terms, and shall remain cancellable by the member that placed it.
  \rationale{The member was told the order was accepted, and has no reason to think a venue
  process has restarted. An order that stops existing without any report is a
  failure the member cannot detect, cannot cancel and cannot hedge against.}
  \covers{ha_test:50}
\end{requirement}

\begin{requirement}{R-0120}{A cancel for an order the venue no longer holds says which case it is}
  Where a member asks to cancel an order the venue does not currently hold, it
  shall be told whether the order once existed and is finished, or was never
  known to the venue at all.
  \rationale{These are different facts about the member's own position and it
  cannot establish either without being told. The protocol distinguishes them for
  that reason.

  There is a second reason particular to this venue. \emph{Unknown order} is also
  what a member is told when the venue has lost what it was holding. Answering both
  cases the same way leaves a member unable to distinguish a venue that has lost
  its orders from one that is answering correctly. Distinguishing the two cases is
  what stops that failure appearing as an ordinary reply.}
  \uncovered{no scenario cancels an order that has already been cancelled}
\end{requirement}

\begin{gap}{}
  A cancel for an order the venue does not currently hold is answered as an order
  it does not recognise, whether the order was cancelled a moment ago or never
  existed. A member cannot tell those apart, and cannot tell either of them from
  the venue having lost what it held.
\end{gap}

\begin{requirement}{R-0116}{The venue recovers what it held before deciding what to do with it}
  A restarted instance shall rebuild the set of orders it held before deciding
  whether to resume trading or to cancel them.
  \rationale{Cancelling an order requires knowing which orders to cancel and which member to
  tell about each one. A venue that cancels without having recovered cannot report
  anything, and stops mentioning the orders instead, which is what \reqref{R-0020}
  exists to prevent. Recovery is therefore required for both outcomes, and it is
  what makes cancelling possible.}
  \covers{ha_test:51}
\end{requirement}

\begin{requirement}{R-0123}{A member keeps what it had, even when the venue's own record of it is unusable}
  Where the record a component keeps of the orders it held cannot be used ---
  absent, damaged, or written by a different build --- the venue shall establish
  what was open from its record of what it took. Where that record reaches back to
  the first command of the trading day, the component shall hold those orders as
  they were and carry on, and shall report at Error that it rebuilt its record.
  Where it does not reach back that far, the venue shall halt trading, through a
  command in its record so that every component halts at the same point, and say
  that it cannot account for what it was holding.
  \rationale{\reqref{R-0102} says such a record is not to be used. That leaves
  the component unable to name what it held --- so without a second route it
  would fall silent about every order, which is the outcome \reqref{R-0020}
  exists to prevent.

  The record of what the venue took is that second route. The component applies
  the same commands in the same order and reaches the same result whatever it did
  before, so replaying that record from the first command of the day gives exactly
  the orders the venue held. Cancelling them would cost every member its orders
  and protect nothing, so the component keeps them. Rebuilding is far slower than
  reading a checkpoint, so it is not the ordinary recovery; this case is rare, and
  it is reported at Error so that someone looks at the record that could not be
  used.

  Where the record does not reach the first command, the orders taken before it
  cannot be named, and cancelling the ones that can be named would leave the
  others unmentioned. Halting and saying that the venue cannot account for what it
  holds is the last resort, and it is still better than silence. The halt is a
  command in the record, so that no component trades on after another has
  stopped.}
  \covers{ha_test:52}
\end{requirement}

\begin{requirement}{R-0117}{Orders open across a long absence are cancelled, reported, and trading halts}
  Where the venue was unable to match for longer than a stated period, it shall
  halt, cancel every order open when it stopped, and send the member that placed
  each a report saying so. Every order it took and had not yet matched shall be
  rejected, with a report saying that trading halted while the order was waiting.
  \rationale{\textbf{The period is the length of the outage, not the age of the
  order.} An order that has rested since the morning in a working market is not a problem,
  because its owner could have cancelled it at any moment and chose not to. The
  hazard is that the owner was \emph{locked out of the order}: the venue could not
  match, so the member could not cancel, and the price moved while the member could
  do nothing. The order is now priced for a market that no longer exists, and its
  owner had no opportunity to act on that.

  The hazard therefore does not depend on how long the order had been resting, nor
  on whether the order was recoverable. An order recovered correctly, whose member
  was locked out of it for ten minutes, carries the same hazard as one recovered
  from a damaged region.

  Cancelling removes the uncertainty. The reports tell each member where it stands.
  The halt records that the venue did not simply carry on, and it does not depend on
  whether any order was open: the venue halts because it could not recover by itself
  within the period it states, and a person decides when trading resumes.

  \textbf{An order taken while the venue could not match is rejected, not matched
  late.} The venue records such an order and holds it, so that it can be matched
  after an ordinary failover, which is short. Once the absence has passed the stated
  period, matching it late would put an order on a book the venue has just cancelled
  for being out of date, and its member has been told nothing about it all the while.
  Rejecting it gives the member a definite answer, and the same answer whatever
  happens next. An order to cancel that was waiting is carried out, because applied
  late it still does what the member asked.

  The venue can measure both parts of this condition, the length of the absence and
  what it cancelled, so it can enter the condition without a person present.
  Lifting the condition still requires a person, under \reqref{R-0023}.}
  \covers{ha_test:51, ha_test:85}
\end{requirement}

\begin{requirement}{R-0118}{How long an absence may be before orders are cancelled is a stated figure}
  The period in \reqref{R-0117} shall be a figure the venue states, measured from
  the moment the venue stopped being able to match to the moment it resumes, and
  the venue shall determine that without depending on anything outside the
  component concerned.
  \rationale{The figure covers the whole absence rather than the recovery step alone.
  Noticing the death, restarting the process and rebuilding the set all count
  towards it, and the restart alone was measured at about two and a half seconds. A
  figure set against recovery time is set against the wrong quantity, and is
  reached later than intended.

  \textbf{It is measured from when matching stopped, not from the last order
  matched.} Traffic comes in bursts, so a quiet period before a failure is ordinary. Taking
  the last order matched as the start would read a ten-minute quiet period followed
  by a two-second restart as a twelve-minute absence, and would cancel every
  member's orders because the market had been quiet.

  \textbf{The component must be able to establish the figure alone.} A supervisor
  records when it killed a process and when it started the replacement, so it is
  the obvious source, but nothing in this chapter may depend on a supervisor being
  present. A component started by hand, in a debugger, hours after the last one
  stopped must reach the same answer, and it does: an hour has passed, the orders
  are an hour stale, and cancelling them is correct.

  \textbf{The figure is the longest recovery the venue makes by itself.} It is set
  from how long an ordinary failover takes, with a margin, so that an ordinary failover
  never reaches it, and every absence longer than that is treated as the venue being
  unable to recover without a person. Where a recovery under way says how long it
  will still take, the venue may act as soon as the absence and that estimate
  together pass the figure, rather than waiting for the absence alone to pass it.}
  \covers{ha_test:51}
\end{requirement}

\begin{requirement}{R-0036}{A cancelled order does not become open again}
  An order the member cancelled before a restart shall not be open after it.
  \rationale{Recovery reads a record of the state and then applies what followed it. A
  cancellation that falls between the two can leave the order open again. The
  member was told the order was cancelled, so it is not watching for that order.}
  \uncovered{no scenario cancels an order and then restarts the engine}
\end{requirement}

\begin{requirement}{R-0020}{An order the venue cannot recover is cancelled to its member}
  Where the set of open orders cannot be rebuilt, the venue shall send an
  unsolicited cancel for every order it can still name, and shall say that it
  cannot name the rest.
  \rationale{A member's view of its orders matches the venue's only if the venue reports what
  it has lost. Ceasing to mention an order reports nothing.

  This requirement is the fallback beneath \reqref{R-0073}, and applies where
  recovery was attempted and could not complete.}
  \uncovered{no scenario makes recovery fail and then checks what the member is sent}
\end{requirement}

\begin{requirement}{R-0021}{A member is not told an order is unknown while the venue is still recovering}
  A request concerning an order shall not be refused as unrecognised before the
  set of open orders has been brought up to date.
  \rationale{Until the set is up to date, the venue cannot establish whether it holds the
  order. Answering that the order is unrecognised makes a statement the venue
  cannot support, and the member acts on that statement.}
  \uncovered{no scenario asserts that open orders survive a restart}
\end{requirement}

\begin{requirement}{R-0108}{The work needed to bring a component current is bounded}
  A checkpoint shall be taken often enough that the records a recovering
  component must apply after reading one do not exceed a stated bound, and that
  bound shall be a figure the venue can state.
  \rationale{Recovery takes as long as reading the checkpoint plus applying everything
  recorded since it. An unstated interval between checkpoints leaves the recovery
  time unstated too, and the interval is then set by whatever was convenient rather
  than by how long the venue may be absent.}
  \uncovered{no scenario measures the records applied after a checkpoint is read, and no bound is stated to measure against}
\end{requirement}

\begin{requirement}{R-0102}{A checkpoint that cannot be trusted is not used}
  A component shall establish that a checkpoint is whole before rebuilding from
  it, and shall not act on the part it could read. A checkpoint that is absent,
  incomplete or inconsistent shall be treated as no checkpoint at all.
  \rationale{A component reads a checkpoint once, at the moment no other record of what it
  holds is available, and it acts on the whole of what it reads. A partial
  checkpoint produces a set of orders that is wrong, and no later check detects the
  error, because every subsequent check is made against that set. Treating a
  doubtful checkpoint as absent gives the outcome of \reqref{R-0020}, which is
  worse for a member but is known.}
  \covers{ha_test:52}
\end{requirement}

\begin{requirement}{R-0022}{The venue keeps enough history to bring a recovered set of orders up to date}
  The record of committed orders shall be retained at least as far back as the
  oldest checkpoint any component may restart from.
  \rationale{Recovery reads a checkpoint and then asks for what followed it. If what followed
  has been discarded, that recovery fails in the way hardest to notice: the
  checkpoint loads, the catch-up returns nothing, and the resulting set appears
  complete.}
  \uncovered{no scenario exhausts retention, which is anchored to nothing measurable}
\end{requirement}

\subsection{Orders that outlive the trading day}
\label{sec:ha:disabled:long_lived}

Most orders last for the session that placed them, or for the trading day. A
member may also place an order to rest for weeks, and expects to find that order
open each morning without doing anything.

Everything else in this section concerns state that must survive a process or a
machine. These orders must survive \emph{the venue}: every process stopped
deliberately, the day ended, the machines rebooted, and trading started again the
next morning.

\subsubsection{Why the mechanisms already described do not serve}
\label{sec:ha:disabled:long_lived:why}

A checkpoint records the open orders at a particular sequence number. It can
be used only if the log still contains every record after that number, so the
venue can apply those records and bring the orders up to date. A day boundary
breaks that relationship, because a new day's record does not continue the
previous day's. Retaining every day's record, so that a checkpoint from any day
remains usable, would require the venue to keep every record it has ever
written.

An order intended to outlive the day therefore needs a durable store of its own.
The venue holds that store independently of the record of what it did, and reads
it when trading begins rather than recovering it when something fails.

\textbf{No component reaches that store while it is trading.} The venue's
components hold no connection to the store, for the same reason that admission
does not depend on one: where the venue needs the store to be reachable in order
to work, every outage of the store becomes a trading outage. The venue exports
the orders and reads them when trading begins. It brings the store up to date
afterwards from its own record of what it did. The store therefore lags the
venue, and the venue can bring it level again from that record.

That has one consequence for the requirement below. \emph{The venue's own
record} must be durable before the venue tells a member its order has been
cancelled. The store is brought up to date afterwards.

\subsubsection{What the venue must do}
\label{sec:ha:disabled:long_lived:behaviour}

\begin{requirement}{R-0105}{An order that outlives the trading day is held where it survives the venue stopping}
  An order whose terms say it remains open beyond the trading day shall be
  recorded where it survives every component stopping and the day ending.
  \rationale{The venue told the member that the order would rest until the member said
  otherwise. Holding the order only where the day's own state is held makes that
  true for one day, which is not what the member was told.}
  \uncovered{no scenario stops the venue and starts it again with such an order open}
\end{requirement}

\begin{requirement}{R-0106}{Such orders are open when trading next begins}
  When trading begins, every order that was open at the end of the previous day
  and has not since expired shall be open, on the terms it was placed under, and
  cancellable by the member that placed it.
  \rationale{A member that must place its resting orders again each morning does not have
  orders that outlive the day. It has orders that the venue cancels overnight
  without reporting the cancellation.}
  \uncovered{no scenario spans a trading day boundary, which the venue does not mark}
\end{requirement}

\begin{requirement}{R-0114}{The store is known to be readable before trading begins, not during it}
  Whether the store holding these orders can be read shall be established before
  trading starts, and its being unreadable shall stop trading from starting
  rather than be discovered once it has.
  \rationale{The venue reads these orders once, at the start of a day. A store that cannot be
  read at that moment leaves the venue opening with an empty set and nothing to
  indicate that anything is missing. Every member holding such an order finds it
  gone. The check costs little, and the venue must make it before it admits
  members.}
  \uncovered{no scenario starts trading with the store unreadable}
\end{requirement}

\begin{requirement}{R-0107}{A change to such an order is committed before the member is told of it}
  Where such an order is cancelled, amended or filled, that shall be committed to
  the venue's own record before the member is told, and the store the order
  outlives the day in shall be brought up to date from that record.
  \rationale{A member told that its order is cancelled stops watching for that order. Where
  the venue writes nothing durable first, a failure between the two brings the
  order back the next time the venue reads the store. The order is then live, on a
  market that has moved, for a member with no reason to look at it. Requiring the
  venue's own record rather than the store is what allows the store to lag. The
  venue can supply the store from that record with anything the store is missing,
  and it can supply nothing for a change that was never committed anywhere.}
  \uncovered{no scenario cancels such an order and then restarts the venue}
\end{requirement}

\begin{gap}{}
  Orders that outlive the trading day are held only where every other open order
  is held, which does not survive the venue stopping. A member placing one is
  given an order that lasts until the next time the venue is restarted, and is
  told nothing to that effect.
\end{gap}

\subsection{The matching engine publisher}
\label{sec:ha:disabled:mep}

The publisher takes what the venue has done and makes it available to consumers
inside the venue: market data, and the recorder described in
\cref{sec:ha:disabled:oar}. The publisher dies and is restarted. No member
notices this directly. What stops is the flow of records to those consumers.

\subsubsection{What state it holds}
\label{sec:ha:disabled:mep:state}

The publisher holds the record of what it has been given, and how far each
subscriber has consumed that record. It does not have to retain the second of
those across a restart. A subscriber presents its own position when it
reconnects, so the publisher is told where each subscriber has reached.

\subsubsection{Where the copy that outlives it lives}
\label{sec:ha:disabled:mep:durable}

The copy lives in the publisher's own log, on disk. The slowest subscriber that
has not yet consumed a record determines what the publisher may discard from that
log.

\begin{gap}{BUG-0048}
  What the publisher may discard is anchored to no position at all. It discards
  nothing, so its log grows for as long as the venue runs, and the positions its
  subscribers have reached are never what decides that a record is kept.
\end{gap}

\subsubsection{How a restarted instance becomes current}
\label{sec:ha:disabled:mep:recovery}

The publisher reopens its log and waits for each returning subscriber to present
its position. The publisher has one such operation rather than two, as the
gateway in \cref{sec:ha:disabled:gateway} does.

\subsubsection{What the member is told}
\label{sec:ha:disabled:mep:member}

Nothing. A member is not a subscriber. The requirements in this subsection
protect the consumers instead.

\begin{requirement}{R-0047}{A subscriber resumes where it had reached, across a restart of the publisher}
  The venue shall supply a returning subscriber with every record after the
  position it presents, and shall not require it to have been present
  throughout.
  \rationale{A subscriber that must stay connected continuously to avoid missing records
  cannot be restarted, upgraded, or briefly separated from the network. No operator
  could run such a subscriber.}
  \uncovered{no scenario restarts the publisher under a connected subscriber}
\end{requirement}

\begin{requirement}{R-0048}{A record is retained until every subscriber has consumed it}
  The venue shall not discard a record that any subscriber entitled to it has
  not yet consumed.
  \rationale{Where the publisher discards records on a timer rather than when they have been
  consumed, a subscriber that was briefly away loses records and nothing reports
  the loss. The subscriber discovers the loss, and the publisher that caused it
  does not.}
  \uncovered{no scenario returns a subscriber after its records were discarded}
\end{requirement}

\subsection{The order activity recorder}
\label{sec:ha:disabled:oar}

The recorder subscribes to the publisher and forwards what the venue has done to
an enterprise bus, which is how systems outside the venue learn of it. A failure
of the recorder is felt outside the venue rather than within it.

\subsubsection{What state it holds}
\label{sec:ha:disabled:oar:state}

The recorder holds one number: how far it has published and had confirmed. It can
obtain everything else again from the publisher.

\subsubsection{Where the copy that outlives it lives}
\label{sec:ha:disabled:oar:durable}

The copy is the recorder's own record of that number, written to disk. The
recorder \textbf{writes it only after the bus has confirmed the publish, and
never before}. That ordering is what makes the design durable. A failure between
publishing and recording the position causes the recorder to publish the record
again, and the consumers downstream handle the duplicate. A failure in the other
order skips a record permanently, and nothing afterwards can detect the omission.

The recorded position therefore needs to be no more than a lower bound. The
recorder can write it as late as it likes, whatever the position's staleness,
because a stale position costs duplicate records and nothing else.

\subsubsection{How a restarted instance becomes current}
\label{sec:ha:disabled:oar:recovery}

The recorder reads its position, presents it to the publisher, and continues. It
publishes again any record it had already published but not recorded. This
duplication is expected: a consumer of the enterprise bus is required to cope
with records it has already seen.

\subsubsection{What the member is told}
\label{sec:ha:disabled:oar:member}

Nothing directly. The last requirement in this subsection is the exception,
because it halts trading.

\begin{requirement}{R-0049}{Nothing is lost from the external record, though it may appear more than once}
  Every record the venue produces shall reach the enterprise bus at least once.
  Duplicates are permitted; omissions are not.
  \rationale{The venue cannot repair a loss outside the venue, because nothing outside it
  knows which records are missing. Anybody holding the record's identity can repair
  a duplicate.}
  \uncovered{no scenario, and none is possible until the recorder exists}
\end{requirement}

\begin{requirement}{R-0050}{Every externally published record carries an identity a consumer can deduplicate on}
  Each record shall carry a stable producer identity and a sequence number
  assigned by the venue.
  \rationale{This requirement is what makes \reqref{R-0049} usable. Without an identity, a
  consumer cannot distinguish a duplicate from a second event, and permitting
  duplicates would then permit the record to be wrong.}
  \uncovered{no scenario, and none is possible until the recorder exists}
\end{requirement}

\begin{requirement}{R-0051}{The recorder never resumes beyond its last confirmed publish}
  The position the recorder resumes from shall never be later than the last
  publish the bus confirmed.
  \rationale{Recording a position before the bus confirms the publish turns a crash into a
  permanent omission, which \reqref{R-0049} forbids. Relaxing the ordering to
  increase throughput would reintroduce that omission.}
  \uncovered{no scenario, and none is possible until the recorder exists}
\end{requirement}

\begin{requirement}{R-0052}{A restart of the recorder does not stop trading}
  The venue shall continue to accept orders while the recorder restarts and
  catches up.
  \rationale{A supervised restart is a routine event, and the venue recovers from it without
  help. Halting trading for such a restart would make the recorder's availability,
  rather than the venue's, decide whether members can trade.}
  \uncovered{no scenario, and none is possible until the recorder exists}
\end{requirement}

\begin{requirement}{R-0053}{Where external publishing cannot resume, trading is halted}
  Where no recorder is publishing and the unpublished backlog has passed a
  stated bound, the venue shall halt trading, and shall say that it has.
  \rationale{Nobody can reconstruct trading afterwards once it has outrun the record of it.
  The venue can measure both parts of this condition: whether a recorder is
  connected, and how large the backlog is. It can therefore act on the condition
  without a person present. The bound must exceed the time a supervised restart and
  catch-up take, or it would contradict \reqref{R-0052}, so the bound has to be
  measured rather than chosen.}
  \uncovered{no scenario, and none is possible until the recorder exists}
\end{requirement}

\subsection{The sequencer and the matching engine together}
\label{sec:ha:disabled:seq_and_me}

Both processes die, and both are restarted. This case is not the two preceding
subsections combined, because each of those recoveries relies on a component that
has also died here.

\subsubsection{What state it holds}
\label{sec:ha:disabled:seq_and_me:state}

Between them, the two components hold everything the venue has about work in
progress. The sequencer holds the record of what has been taken. The matching
engine holds which orders are open. Neither holds the other's state, and neither
can reconstruct it.

\subsubsection{Where the copy that outlives it lives}
\label{sec:ha:disabled:seq_and_me:durable}

Two records survive. The sequencer's log holds what the venue took and the order
it took them in. The matching engine's region holds the orders it had open and
the position those orders are current to, which
\cref{sec:ha:disabled:me:durable} describes.

What does not survive is the sequencer's record that particular orders were
waiting for an engine. Only the running process holds that.

This case differs from the single failures because the component each of those
recoveries relies on has also died. A matching engine becomes current by asking
the sequencer, so it cannot become current until the sequencer is current
itself.

\subsubsection{How a restarted instance becomes current}
\label{sec:ha:disabled:seq_and_me:recovery}

This subsection is about an order that sits between the two components:
committed by the sequencer, and never given to the matching engine.

The engine establishes the position its region is current to, asks the sequencer
for everything committed after that position, and applies what comes back. The
order is in the log, so the sequencer serves it. The engine applies it and sends
the member an execution report for it, marked as a report the member may already
have seen, which is \reqref{R-0122}.

The sequencer's own record that the order was waiting does not survive the
restart, and the venue does not depend on it. The position the engine presents is
what determines the records the sequencer sends.

Recovery now has a required order: the sequencer must become current before the
matching engine can. The engine gets missing records from the sequencer, so it
cannot catch up until the sequencer is ready. This dependency does not arise
when only one component fails, because the other remains running.

\subsubsection{What the member is told}
\label{sec:ha:disabled:seq_and_me:member}

\begin{requirement}{R-0015}{Trading does not resume while committed orders remain unresolved}
  The venue shall not accept new orders after such a failure until every order
  it had committed has either been given to the matching engine or had an
  execution report produced for it. Whether a member is present to receive that
  report does not bear on this: a member that is slow to return, or that never
  returns, cannot hold the venue shut.
  \rationale{Accepting new orders while earlier orders remain unresolved conceals the loss
  rather than recovering from it. The venue then looks healthy to a member whose
  earlier order the venue will never mention again.}
  \uncovered{no scenario restarts both components together}
\end{requirement}

Where the automatic paths cannot resolve a failure, a person resolves it, and a
longer timer does not. The venue therefore needs a halted state that a person can
put it into. The venue also needs that state for reasons unconnected to any
failure: a scheduled pause, a market-wide event, or an instrument trading in
disorderly conditions.

\begin{requirement}{R-0023}{Trading can be halted by a person, and members are told}
  An operator shall be able to halt trading. Members shall be told that the
  venue has halted and shall be told when it resumes, and the halt shall not
  lift of its own accord.
  \rationale{A member that is connected but idle sends no order, so the rejection of an order
  tells it nothing. A member reconnecting into a halt would otherwise discover the
  halt by sending an order and having it refused. Requiring a person to lift the
  halt is what distinguishes it from the automatic states that clear themselves.}
  \covers{ha_test:94, ha_test:95}
\end{requirement}

\begin{gap}{BUG-0065}
  An operator halts trading and lifts the halt, and a member placing an order
  during a halt is told in the reply that the venue is halted. A member that is
  connected and not trading, or that logs on during a halt, is told nothing, and
  no member is told when trading resumes. \reqref{R-0023} is met only in part
  until then.
\end{gap}

\begin{requirement}{R-0024}{A halt that cancels open orders reports every cancellation}
  Where a halt cancels the venue's open orders, an execution report shall be
  sent to the member that placed each one. No order shall cease to be open
  without the member that placed it being told.
  \rationale{Cancelling the open orders is a normal thing for a halt to do, and real venues
  have done it when a market became disorderly. What makes it safe is that each
  member's view of its orders comes to match the venue's, and that requires the
  venue to report what it cancelled rather than to stop mentioning those orders.}
  \covers{ha_test:51, ha_test:85}
\end{requirement}

\begin{requirement}{R-0017}{Recovery does not depend on the order the components return in}
  The venue shall reach the same state whichever of the sequencer and the
  matching engine starts first after both have failed.
  \rationale{A launcher restarts what it starts and coordinates with nothing, so the venue
  does not decide which of the two components returns first. A component that
  cannot yet be told what it needs must wait, rather than proceed on the assumption
  that there is nothing to receive.}
  \uncovered{no scenario restarts both components together}
\end{requirement}

\subsection{The gateway together with another component}
\label{sec:ha:disabled:gw_and_others}

The gateway depends on the sequencer for everything it recovers, and on the
matching engine for everything it asks to be done on a member's behalf. Losing
either of those alongside the gateway is therefore a separate case. This
subsection states each case even where the requirements are stated elsewhere, so
that a reader can see which combinations the venue is designed to survive.

\subsubsection{The gateway and the sequencer}
\label{sec:ha:disabled:gw_and_seq}

A member reconnects to a restarted gateway, and that gateway asks a sequencer
that has itself just restarted and no longer holds the session. This case needs
no new requirement. \reqref{R-0011} requires the numbering to survive, and
\reqref{R-0012} requires the venue to say so where it has not, rather than
resuming the member at the beginning. The combination is the likeliest route to
\reqref{R-0012}, because no component reports a session that was not bound at the
moment the sequencer restarted.

\subsubsection{The gateway and the matching engine}
\label{sec:ha:disabled:gw_and_me}

A member's connection ends while the matching engine is not running. The gateway
does not act on the member's orders itself. It reports the disconnection and those instructions to the component that holds
the orders, and that component is absent.

\begin{requirement}{R-0037}{A member's cancel-on-disconnect instruction is not lost because a component was restarting}
  Where a session ends while the component that holds its orders is unavailable,
  the instruction shall be acted on once that component returns, or the member
  shall be told that it was not.
  \rationale{The member configured this to bound its exposure when its own connection fails.
  An instruction dropped because a venue process was restarting leaves the member
  exposed, and neither the member nor the venue holds any record of it.}
  \uncovered{no scenario ends a session while the engine is down}
\end{requirement}

\subsection{The member's own application}
\label{sec:ha:disabled:client}

A member's application dies and is restarted. This is the only failure in this
section that happens outside the venue, so the venue's obligations differ. The
venue has nothing to recover here, and must hold everything it was keeping for
that member.

It is also the failure the venue meets most often. A member restarts its systems
on its own schedule, for its own reasons, and tells the venue nothing beforehand.

\subsubsection{What state it holds}
\label{sec:ha:disabled:client:state}

The member's application holds its own view: which orders it treats as live at
the venue, the identifiers it has used, and how far its side of the session
numbering had reached. The venue does not keep any of this, and must not assume
that any of it survived.

The venue holds a different set of state \emph{on the member's behalf}, and the
rest of this subsection concerns it: the session's numbering, the orders open
under that session, and the terms the member was admitted under.

\subsubsection{Where the copy that outlives it lives}
\label{sec:ha:disabled:client:durable}

The member chooses where to keep its own copy. The venue holds no record of that
choice and must not depend on it. A member may return holding its full position,
or holding none of it, and both are legitimate.

The venue must therefore answer the member rather than assume anything about it.
A returning member that has kept its numbering resumes at that numbering. A
member that has lost its numbering says so and starts again from the beginning.
In both cases the member must be able to establish what the venue holds for it,
which \cref{sec:ha:disabled:gateway:member} records that it currently cannot.

\subsubsection{How a restarted instance becomes current}
\label{sec:ha:disabled:client:recovery}

The member's application connects, logs on, and either resumes its numbering or
asks to restart it. The venue supplies what it holds for the session, and the
member reconciles its own view against what the venue reports.

\subsubsection{What the member is told}
\label{sec:ha:disabled:client:member}

\begin{requirement}{R-0054}{A member that returns in time finds its orders untouched}
  Where a member's connection ends unexpectedly and the same comp id logs on
  again within the period provisioned for it, none of that member's orders shall
  have been cancelled.
  \rationale{The period exists so that a member's own restart is a reconnection rather than an
  event the market sees. Cancelling the orders and having the member place them
  again costs the member its position in the queue, and reports activity to the
  market that did not happen.}
  \covers{ha_test:19}
\end{requirement}

\begin{requirement}{R-0055}{A member that does not return has its orders cancelled as it asked, and is told when it comes back}
  Where the period expires without the member returning, the orders it asked to
  have cancelled shall be cancelled, and the reports of those cancellations
  shall be delivered to it when it next connects.
  \rationale{The member configured this to bound its exposure while it is absent. Cancelling
  the orders without keeping the reports for the member's return leaves the member
  to work out what happened for itself.}
  \uncovered{no scenario lets the period expire and then reconnects}
\end{requirement}

\begin{requirement}{R-0056}{An order meant to outlive the session is not cancelled by a disconnection}
  An order whose time in force says it survives the trading session shall not be
  cancelled because the member's connection ended.
  \rationale{The member placed such an order to outlive the session that placed it.
  Cancelling it because a connection closed does the opposite of what the member
  asked for, and the member cannot restate the instruction cheaply.}
  \uncovered{no scenario places an order that outlives the session and then drops it}
\end{requirement}

\begin{requirement}{R-0057}{A member that has lost its own numbering can start afresh without losing its orders}
  A member that asks to restart its numbering shall be given a session numbered
  from the beginning, and the orders it holds at the venue shall be unaffected.
  \rationale{Losing a sequence store is an ordinary consequence of a member's own restart, and
  it says nothing about that member's orders. Tying the two together would make a
  member choose between a session it can use and the orders it holds.}
  \uncovered{no scenario restarts a member's numbering with orders open}
\end{requirement}

\subsection{Supervision}
\label{sec:ha:disabled:supervision}

A launcher answers one question: is this component running, and if not, start it.
It takes no part in deciding roles, and nothing in this chapter depends on it
being present. This subsection states the bounds on what it must do.

\begin{requirement}{R-0033}{A component asked to stop is not restarted}
  Where a component is stopped deliberately, whatever supervises it shall not
  start it again.
  \rationale{Without this requirement, nobody can stop the venue. Every stop becomes a
  restart, and an operator has no way to take a component out of service.}
  \uncovered{no scenario stops a supervised component deliberately}
\end{requirement}

\begin{requirement}{R-0038}{A component that cannot start at all is reported at once}
  Where a component cannot be started --- its binary absent, its configuration
  unreadable, a port it needs already held --- whatever supervises it shall
  report the failure immediately and shall not retry as though the attempt might
  succeed.
  \rationale{Retrying a start that cannot succeed turns a fault an operator could fix in a
  minute into an absence that nothing reports. None of these causes is transient,
  so waiting achieves nothing. A component that starts and then dies is a different
  case, and \reqref{R-0039} covers it.}
  \uncovered{no scenario starts a component that cannot start}
\end{requirement}

\begin{requirement}{R-0039}{A component that fails repeatedly and quickly is not restarted again}
  Where a component dies repeatedly within a short time of being started,
  whatever supervises it shall stop restarting it. A death soon after a start is
  a failure to start rather than a run that ended, and shall be counted as one.
  \rationale{Restarting a component whose fault is deterministic produces another failure, and
  restarting it without limit produces a fault that runs continuously and reports
  nothing. Distinguishing a failure to start from a run that ended is what makes
  the condition detectable: a healthy component runs for hours, and one that exits
  in seconds has not run.}
  \uncovered{no scenario makes a component fail repeatedly}
\end{requirement}

\begin{requirement}{R-0040}{Giving up is announced where something other than a person can see it}
  Having stopped restarting a component, the supervisor shall say so on its
  output, in a form a log monitor can detect and raise an alert from without
  anyone watching at the time.
  \rationale{Automatic restart removes the symptom that would otherwise draw attention.
  Somebody downstream eventually notices a component that dies and stays dead.
  Nobody notices a component that is restarted forty times an hour. Automating the
  recovery therefore makes the reporting necessary.}
  \uncovered{nothing asserts what the supervisor says when it gives up}
\end{requirement}

\begin{requirement}{R-0041}{A component that runs healthily again is given a fresh allowance}
  Where a restarted component runs for long enough to have started
  successfully, the count of rapid failures against it shall be reset.
  \rationale{Without this requirement, the allowance is consumed over the life of the
  deployment rather than over one episode. A component that fails once a month
  eventually exhausts the allowance, and the supervisor then stops restarting it
  for a fault it has recovered from every time.}
  \uncovered{no scenario makes a component fail repeatedly}
\end{requirement}

\subsection{Stopping and starting deliberately}
\label{sec:ha:disabled:stopping}

Not every stop is a failure. An operator takes a component down for a release,
for a configuration change, or because the trading day has ended. None of those
should cost the venue what a crash costs.

\begin{requirement}{R-0034}{An orderly stop leaves less to recover than a crash}
  A component asked to stop shall record its state before it goes, so that what
  it has to recover afterwards is no more than what it wrote.
  \rationale{A planned restart is the common case, for a release or a configuration change,
  and it need not cost what an unplanned restart does. It is also the cheapest way
  to exercise the recovery path regularly.}
  \uncovered{no scenario stops a component deliberately, and nothing would be written if one did}
\end{requirement}

\begin{requirement}{R-0035}{The start and end of a trading day are published}
  The venue shall publish the start of a trading day and its end as events that
  any interested party may subscribe to.
  \rationale{A venue that has just started is either resuming a trading day or beginning a new
  one, and which of those it is changes what every recovery in this section means.
  A component must be able to ask rather than infer it. Publishing the event once
  to every subscriber also stops each component being told separately, which is how
  components come to hold different answers.}
  \uncovered{no scenario spans a trading day boundary, which the venue does not mark}
\end{requirement}

